\section{CFD-PCO Observer Design}

The core idea of CFD-PCO is to perform one-step-ahead velocity prediction using the CFD-pre-calibrated dynamics model, map the prediction residual to a current estimate gradient via sensitivity analysis, and update only when the excitation condition is satisfied.

\subsection{Velocity Prediction Architecture}

The observer predicts the velocity at step $k+1$ using the current estimate $\hat{\bm{c}}_k$:

\begin{equation}
    \hat{\bm{\nu}}_{k+1|k} = \bm{\nu}_k + \Delta t \cdot \bm{a}_{\text{model}}(\bm{\nu}_k, \bm{\tau}_k, \hat{\bm{c}}_k)
    \label{eq:prediction}
\end{equation}

where $\bm{a}_{\text{model}}$ is the acceleration computed from the CFD dynamics model. In contrast to acceleration-residual approaches, velocity prediction leverages the smoothing effect of time integration, avoiding direct measurement or differentiation of $\dot{\bm{\nu}}$---the latter yields a signal-to-noise ratio of only approximately 0.03 under typical DVL noise (0.01\,m/s) at 100\,Hz sampling, rendering gradient-based estimation impractical.

\subsection{Sensitivity Analysis}

Define the sensitivity matrix of the velocity prediction with respect to the current parameters:

\begin{equation}
    \bm{\Phi}_c = \frac{\partial \hat{\bm{\nu}}_{k+1|k}}{\partial \bm{c}} \in \R^{6\times 2}
    \label{eq:sensitivity}
\end{equation}

Although the CFD drag model $\bm{D}(\bm{\nu}_r)$ admits closed-form partial derivatives (being quadratic in $\bm{\nu}_r$), the full 6-DOF analytical differentiation is cumbersome. We compute $\bm{\Phi}_c$ via numerical perturbation:

\begin{equation}
    \bm{\Phi}_c(:,j) = \frac{\hat{\bm{\nu}}_{k+1|k}(\hat{\bm{c}}+\delta\bm{e}_j) - \hat{\bm{\nu}}_{k+1|k}(\hat{\bm{c}})}{\delta}, \quad j=1,2
    \label{eq:numerical_phi}
\end{equation}

with perturbation step $\delta=0.01$\,m/s and standard basis vectors $\bm{e}_j$. Each update requires 3 CFD model evaluations (1 baseline + 2 perturbations), which is computationally tractable.

The sensitivity Gramian is defined as:

\begin{equation}
    \bm{W}_c = \bm{\Phi}_c^T\bm{\Phi}_c \in \R^{2\times 2}
    \label{eq:gramian}
\end{equation}

$\bm{W}_c$ approximates the Fisher information matrix for current parameter estimation. Its minimum eigenvalue $\lambda_{\min}(\bm{W}_c)$ quantitatively characterizes the identifiability of current parameters under the current dynamic state.

\subsection{Continuous Gauss-Markov Regularization}

Rather than employing a binary excitation gate that switches between gradient updates and GM propagation, PC-RCO applies the GM decay \emph{continuously} at every time step, in parallel with the gradient update:

\begin{equation}
    \hat{\bm{c}}_{k+1} = e^{-\Delta t/\tau_c}\left(\hat{\bm{c}}_k + \Delta\bm{c}_k^{\text{grad}}\right) + (1-e^{-\Delta t/\tau_c})\bar{\bm{c}}
    \label{eq:continuous_gm}
\end{equation}

where $\Delta\bm{c}_k^{\text{grad}}$ is the gradient-driven update (Section~3.4) and $\bar{\bm{c}}$ is the prior mean. This continuous formulation has two regimes:
\begin{itemize}[itemsep=0pt,topsep=4pt]
    \item \textbf{High excitation}: $\|\Delta\bm{c}_k^{\text{grad}}\| \gg (1-e^{-\Delta t/\tau_c})\|\hat{\bm{c}}_k - \bar{\bm{c}}\|$ --- the gradient update dominates, and the GM decay acts as weak Tikhonov regularization.
    \item \textbf{Low excitation}: $\|\Delta\bm{c}_k^{\text{grad}}\| \ll (1-e^{-\Delta t/\tau_c})\|\hat{\bm{c}}_k - \bar{\bm{c}}\|$ --- the GM decay dominates, gradually pulling the estimate toward the prior and preventing drift.
\end{itemize}

\textbf{Motivation for continuous over binary gating}: Through systematic diagnosis, we found that binary excitation gates based on sensitivity Gramian eigenvalues ($\lambda_{\min}(\bm{W}_c) > \mu_{\text{gate}}$) remain permanently open in feedback-controlled AUVs. The root cause is twofold: (1) the CFD drag model's quadratic form provides nonzero sensitivity even at very low turn rates; (2) feedback controllers (e.g., ALOS guidance) continuously generate small heading corrections to counteract cross-track drift, preventing sustained zero-excitation conditions. The continuous approach eliminates threshold tuning and provides robust regularization regardless of excitation level.

\subsection{Gradient Update with Adaptive Clamping}

The gradient-driven update $\Delta\bm{c}_k^{\text{grad}}$ in Eq.~\eqref{eq:continuous_gm} is computed as:

\begin{equation}
    \Delta\bm{c}_k^{\text{grad}} = \gamma \cdot \frac{\bm{\Phi}_c^T \bm{e}_{\nu}}{1 + \lambda_{\min}(\bm{W}_c) \cdot \kappa}
    \label{eq:update}
\end{equation}

where $\bm{e}_{\nu} = \bm{\nu}_{k+1}^{\text{meas}} - \hat{\bm{\nu}}_{k+1|k}$ is the velocity prediction innovation (surge and sway components only), $\gamma = K_{\text{obs}} \cdot 100$ is the effective gain, and $\kappa$ is the adaptive regularization coefficient. We choose the gradient direction $\bm{\Phi}_c^T\bm{e}_{\nu}$ over the Newton direction to avoid numerical instability when the Gramian is near-singular and to provide greater robustness under model mismatch.

\textbf{Adaptive clamping with online noise estimation}: To balance convergence speed against physical plausibility, we apply per-step clipping with a dynamically adjusted bound:

\begin{equation}
    \|\Delta\bm{c}_k^{\text{grad}}\| \leq \Delta c_{\max}(\hat{\sigma}_k), \quad \Delta c_{\max}(\hat{\sigma}_k) = \min\left(\Delta c_{\max}^{\text{base}} \cdot \max\left(0.5,\; \min\left(3.0,\; \frac{\hat{\sigma}_k}{\sigma_{\text{ref}}}\right)\right),\; 0.20\right)
    \label{eq:adaptive_clamp}
\end{equation}

where $\hat{\sigma}_k$ is an exponential moving average (EMA) estimate of the instantaneous DVL noise level computed from the velocity prediction residual:

\begin{equation}
    \hat{\sigma}_k = 0.995\,\hat{\sigma}_{k-1} + 0.005\,\|\bm{e}_{\nu}\|
    \label{eq:noise_est}
\end{equation}

$\sigma_{\text{ref}}=0.02$\,m/s is the reference DVL noise level, and $\Delta c_{\max}^{\text{base}}=0.06$\,m/s is the nominal clamp bound. The adaptive scheme: (1) maintains a tight clamp ($\approx$0.06\,m/s) under normal low-noise conditions for smooth convergence; (2) relaxes to approximately 0.18\,m/s under high noise to maintain convergence through noisy measurements; (3) enforces an absolute upper bound of 0.20\,m/s to prevent unbounded updates under any conditions.

\subsection{Stability Analysis (Sketch)}

Define the Lyapunov candidate:

\begin{equation}
    V = \frac{1}{2}\tilde{\bm{\nu}}^T\bm{P}\tilde{\bm{\nu}} + \frac{1}{2\gamma_c}\|\tilde{\bm{c}}\|^2
    \label{eq:lyapunov}
\end{equation}

where $\tilde{\bm{\nu}} = \bm{\nu} - \hat{\bm{\nu}}_{obs}$ is the state estimation error and $\tilde{\bm{c}} = \bm{c} - \hat{\bm{c}}$ is the current estimation error. Through analysis of $\dot{V}$, it can be shown that under CFD model error $\Delta\bm{D}$, the tracking error is uniformly ultimately bounded (UUB), with the error bound proportional to the uncertainty upper bound $\bar{\Delta}_D$. The complete Lyapunov derivation is provided in the Appendix.

\subsection{Feedforward Compensation Strategy}

The current estimate $\hat{\bm{c}}$ is ultimately used for feedforward compensation. Converting $\hat{\bm{c}}$ to body-frame current components, the net force/moment induced by the current is:

\begin{equation}
    \hat{\bm{d}} = \bm{M}\left[\dot{\bm{\nu}}(\hat{\bm{c}}) - \dot{\bm{\nu}}(\bm{0})\right]
    \label{eq:feedforward}
\end{equation}

\textbf{Compensation ablation discovery}: During tests injecting individual components of $\hat{\bm{d}}$ into the SMC control law, we discovered an important coupling phenomenon---surge feedforward compensation (\texttt{hat\_d(1)}) severely degrades heading control, worsening heading RMSE from 0.14\,rad to 0.93\,rad (surge-only compensation). The root cause is that main thruster RPM changes indirectly affect yaw moment through the thrust allocation matrix in the XHY AUV's 5-thruster configuration, creating a surge-yaw dynamic coupling. Based on this finding, CFD-PCO employs a \textbf{yaw-only feedforward} strategy: only $\hat{\bm{d}}(6)$ (yaw moment compensation) is injected into the control law, while the surge channel is handled independently by the speed PID controller. This design choice, validated as optimal through ablation experiments, is motivated by engineering practicality and avoids cross-channel coupling.
